% \iffalse meta-comment
%
% Copyright (C) 2019-2026 by Jan Hajer
% -----------------------------------
%
% This file may be distributed and/or modified under the
% conditions of the LaTeX Project Public License, either version 1.3c
% of this license or (at your option) any later version.
% The latest version of this license is in:
%
% http://www.latex-project.org/lppl.txt
%
% and version 1.3c or later is part of all distributions of LaTeX
% version 2005/12/01 or later.
%
% \fi
%
% \iffalse

%<package>\NeedsTeXFormat{LaTeX2e}[2021-11-15]
%<package>\ProvidesPackage{hep-text}[2026/11/13 v1.6 Text utilities: language setup, abbreviations, links, breakable underlining and compact lists]
%<documentation>\ProvidesFile{hep-text-documentation.tex}[2026/11/13 v1.6 HEP-Text documentation]
%<documentation>\ProvidesFile{hep-text-text.tex}[2026/11/13 v1.6 HEP-Text tests]
%
%<*documentation>

\RequirePackage[l2tabu, orthodox]{nag}
\documentclass{ltxdoc}

\renewcommand*\theCodelineNo{\rmfamily\tstyle\footnotesize\arabic{CodelineNo}}
\AtBeginEnvironment{macrocode}{\renewcommand*{\ttdefault}{clmt}}
\renewcommand*{\MacroFont}{\codestyle}
\AtBeginDocument{\DeleteShortVerb{\|}}
\AtBeginDocument{\MakeShortVerb{\"}}
\EnableCrossrefs
\CodelineIndex
\RecordChanges

\usepackage{hologo}

\usepackage[parskip,oldstyle,font=10pt]{hep-paper}
\bibliography{bibliography}

\acronym{CM}{computer modern}
\acronym{LM}{latin modern}
\shortacronym{PDF}{portable document format}

\newenvironment{columns}[1][.5]{%
  \par\vspace{-\bigskipamount}%
  \begin{minipage}[t]{\linewidth}%
  \begin{minipage}[t]{#1\linewidth}%
  \def\column{%
    \end{minipage}%
    \begin{minipage}[t]{\linewidth-#1\linewidth}%
  }%
}{\end{minipage}\end{minipage}\par}
%</documentation>

%<*driver>
\expandafter\newif\csname ifshort\endcsname
\shortfalse
\begin{document}
\DocInput{hep-text-implementation.dtx}
\end{document}
%</driver>
%
% \fi
%
% \CheckSum{406}
%
% \CharacterTable
%  {Upper-case    \A\B\C\D\E\F\G\H\I\J\K\L\M\N\O\P\Q\R\S\T\U\V\W\X\Y\Z
%   Lower-case    \a\b\c\d\e\f\g\h\i\j\k\l\m\n\o\p\q\r\s\t\u\v\w\x\y\z
%   Digits        \0\1\2\3\4\5\6\7\8\9
%   Exclamation   \!     Double quote  \"     Hash (number) \#
%   Dollar        \$     Percent       \%     Ampersand     \&
%   Acute accent  \'     Left paren    \(     Right paren   \)
%   Asterisk      \*     Plus          \+     Comma         \,
%   Minus         \-     Point         \.     Solidus       \/
%   Colon         \:     Semicolon     \;     Less than     \<
%   Equals        \=     Greater than  \>     Question mark \?
%   Commercial at \@     Left bracket  \[     Backslash     \\
%   Right bracket \]     Circumflex    \^     Underscore    \_
%   Grave accent  \`     Left brace    \{     Vertical bar  \|
%   Right brace   \}     Tilde         \~}
%
% \changes{v1.0}{2021/09/01}{First release on CTAN: language and quotation support, foreign abbreviations, text helpers, list environments, and line-breakable underlining.}
% \changes{v1.1}{2022/11/01}{Disabled styling of foreign abbreviations by default, made software names safe in PDF strings, and added \cs{subsubparagraph}.}
% \changes{v1.2}{2023/07/01}{Refreshed the documentation.}
% \changes{v1.3}{2024/11/01}{Switched line-breakable underlining from soulutf8 to soul.}
% \changes{v1.4}{2025/09/01}{Applied starred command definitions consistently.}
% \changes{v1.5}{2026/08/01}{Raised the LaTeX format requirement, suppressed the indentation after lists, and expanded the documentation and README.}
% \changes{v1.6}{2026/11/13}{Restored underline in math mode, made underline, online links, email addresses, and foreign abbreviations safe in PDF strings, fixed the empty-title test of enumdescript, and stopped enumdescript from changing the items of nested lists.}
%
% \ifshort
%<*documentation>
% \fi
%
\GetFileInfo{hep-text.sty}

\title{The \software{hep-text} package\thanks{This document corresponds to \software{hep-text}~\fileversion.}}
\subtitle{Text utilities: language setup, abbreviations, links, breakable underlining and compact lists}
\author{Jan Hajer \email{jan.hajer@tecnico.ulisboa.pt}}
\date{\filedate}

% \ifshort
\begin{document}
% \fi

\newgeometry{vscale=.8, vmarginratio=3:4, includeheadfoot, left=11em, marginparwidth=4.6cm, marginparsep=3mm, right=7em}

\maketitle

\begin{abstract}
The \software{hep-text} package collects lightweight text utilities.
It configures language-aware hyphenation, quotation marks, and Latin abbreviations through established \hologo{LaTeX} packages, adds helpers for software names, online references, email addresses, numero labels, prefix constructions, and line-breakable underlining, and defines compact inline and enumerated description lists based on \software{enumitem}.
\end{abstract}

Load the package with "\usepackage{hep-text}".

\DescribeMacro{lang}
The "lang" option selects the document language.
It accepts the language names understood by \software{babel} \cite{babel}, which is loaded for hyphenation support.

\software{csquotes} \cite{csquotes} provides context-sensitive quotation commands.
\DescribeMacro{\enquote}
\DescribeMacro{\MakeOuterQuote}
The macros "\enquote"\marg{text} and "\MakeOuterQuote{}" leave the choice of quotation marks to \hologo{LaTeX}.
The latter also allows \verb|"| to be used instead of the pair \verb|``| and \verb|''|.

\DescribeMacro{\eg}
\DescribeMacro{\ie}
\DescribeMacro{\cf}
\DescribeMacro{\vs}
The \software{foreign} package \cite{foreign} defines common Latin abbreviations such as "\eg", "\ie", "\cf", and "\vs".
They are typeset as \eg, \ie, \cf, and \vs with appropriate spacing.
Redefine "\foreignabbrfont" to adjust their style, for example "\renewcommand\foreignabbrfont{\itshape}" for italic abbreviations.

\DescribeMacro{\no}
The "\no"\marg{number} macro prints a numero sign followed by a non-breaking space, as in \no{123}.

\DescribeMacro{\software}
The "\software"\oarg{version}\marg{name} macro typesets software names in small caps, with optional version information, as in \software[\fileversion]{HEP-Paper}.

\DescribeMacro{\online}
The "\online"\marg{url}\marg{text} macro links \meta{text} to \meta{url} when \software{hyperref} \cite{hyperref} is loaded and otherwise prints the text in a URL-like style, as in \online{https://ctan.org/pkg/hep-text}{ctan.org/pkg/hep-text}.

\DescribeMacro{\email}
The "\email"\marg{address} macro is shorthand for "\online"\marg{mailto:address}\marg{address}, as in \email{jan.hajer@tecnico.ulisboa.pt}.

\DescribeMacro{\prefix}
The "\prefix"\marg{prefix}\marg{word} macro typesets prefix constructions with a permissible line break after the prefix, as in \prefix{pre}{selection}.

\DescribeMacro{\subsubparagraph}
The "\subsubparagraph" macro starts an unnumbered paragraph-level break with the same vertical spacing as "\paragraph" but without a heading.

\DescribeMacro{inlinelist}
\DescribeMacro{enumdescript}
The "inlinelist" and "enumdescript" environments are built with \software{enumitem} \cite{enumitem}.
\\
\begin{columns}
\begin{verbatim}
The three main points are
\begin{inlinelist}
  \item one
  \item two
  \item three
\end{inlinelist}
\end{verbatim}
\column
% \vspace{.1ex}
The three main points are
\begin{inlinelist}
 \item one
 \item two
 \item three
\end{inlinelist}
\end{columns}
\vspace{4ex}
\begin{columns}[.6]
\begin{verbatim}
\begin{enumdescript}[label=\Roman*)]
  \item{First} one
  \item{Second} two
  \item{Third} three
\end{enumdescript}
\end{verbatim}
\column
\begin{enumdescript}[label=\Roman*)]
 \item{First} one
 \item{Second} two
 \item{Third} three
\end{enumdescript}
\end{columns}

In "enumdescript", "\item"\oarg{label}\marg{title} prints the title in bold; lists nested inside keep their usual "\item".

\DescribeMacro{\underline}
The "\underline" macro is redefined with \software{soul} \cite{soul} so that underlined text can break across lines.
In math mode it keeps its standard meaning.
Text containing macros unknown to \software{soul}, such as "\ref"\marg{key}, "\enquote"\marg{text}, or "\eg", is underlined as a whole without line breaks.

If \software{hyperref} is loaded, software names, online links, email addresses, underlined text, and the foreign abbreviations reduce to their plain text in \PDF strings such as bookmarks.

% \ifshort

\printbibliography

\end{document}
%
%</documentation>
% \fi
%
% \StopEventually{
% \PrintChanges
% \printbibliography
% }
%
% \appendix
%
% \section{Implementation}
%
%<*package>
%
% Load \software{kvoptions} \cite{kvoptions} and set up the "heptext" option namespace.
%    \begin{macrocode}
\RequirePackage{kvoptions}
\SetupKeyvalOptions{
  family=heptext,
  prefix=heptext@
}
%    \end{macrocode}
%
% \begin{macro}{lang}
% Define the "lang" option, which accepts the language names provided by \software{babel} \cite{babel}.
% Use "british" as the default.
%    \begin{macrocode}
\DeclareStringOption[british]{lang}
%    \end{macrocode}
% \end{macro}
%
%    \begin{macrocode}
\ProcessKeyvalOptions*
%    \end{macrocode}
%
% Load \software{babel} \cite{babel} for hyphenation and \software{csquotes} \cite{csquotes} for quotation commands.
%    \begin{macrocode}
\PassOptionsToPackage{\heptext@lang}{babel}
\RequirePackage{babel}
\PassOptionsToPackage{autostyle}{csquotes}
\RequirePackage{csquotes}
%    \end{macrocode}
%
% \begin{macro}{\underline}
% Load \software{soul} \cite{soul} for underlined text that can break across lines.
% Redefine "\underline" robustly so that it uses "\textul" in text mode and keeps its original definition in math mode.
%    \begin{macrocode}
\RequirePackage{soul}
\NewCommandCopy\hep@underline\underline
\DeclareRobustCommand*\underline{%
  \ifmmode
    \expandafter\hep@underline
  \else
    \expandafter\hep@text@underline
  \fi
}
%    \end{macrocode}
% \end{macro}
%
% \begin{macro}{\hep@soul@if@safe}
% Test whether the \software{soul} package can process a text.
% This is the case if every macro in the text, outside of inline math, is known to \software{soul}: registered with "\soulregister", such as font commands and accents, or handled by \software{soul} itself, such as "\\" and "\mbox".
% Other macros may take arguments, which \software{soul} would separate from the macro, or look ahead, as "\xspace" does.
% The same test is defined by the \software{hep-math} package for "\overline", so it is only defined if it does not exist yet.
% "\hep@soul@if@safe" is the \hologo{LaTeXe} name of the conditional "\hep_soul_if_safe:nTF".
%    \begin{macrocode}
\ExplSyntaxOn
\cs_if_exist:NF \hep_soul_if_safe:nTF {
  \seq_new:N \l__hep_soul_seq
  \tl_new:N \l__hep_soul_tl
  \tl_new:N \l__hep_soul_known_tl
  \bool_new:N \l__hep_soul_bool
  \regex_const:Nn \c__hep_soul_math_regex { \cM\$ .*? \cM\$ }
  \regex_const:Nn \c__hep_soul_cs_regex { \c{.*} }
  \tl_const:Nn \c__hep_soul_known_tl {
    \\0\\ \\0\slash \\0\mbox \\0\hbox \\0\soulomit
    \\0\break \\0\linebreak \\0\- \\0\par
  }
  \prg_new_protected_conditional:Npnn \hep_soul_if_safe:n #1 { TF } {
    \group_begin:
    \cs_set:Npn \\ ##1 ##2 { \token_to_str:N ##2 , }
    \tl_set:No \l__hep_soul_known_tl { \the \SOUL@cmds }
    \tl_set:Ne \l__hep_soul_known_tl
      { , \l__hep_soul_known_tl \c__hep_soul_known_tl }
    \tl_set:Nn \l__hep_soul_tl {#1}
    \regex_replace_all:NnN \c__hep_soul_math_regex { } \l__hep_soul_tl
    \regex_extract_all:NVN \c__hep_soul_cs_regex
      \l__hep_soul_tl \l__hep_soul_seq
    \bool_set_true:N \l__hep_soul_bool
    \seq_map_inline:Nn \l__hep_soul_seq {
      \token_if_macro:NT ##1 {
        \exp_args:NNe \tl_if_in:NnF \l__hep_soul_known_tl
          { , \token_to_str:N ##1 , }
          { \bool_set_false:N \l__hep_soul_bool \seq_map_break: }
      }
    }
    \bool_if:NTF \l__hep_soul_bool
      { \group_end: \prg_return_true: }
      { \group_end: \prg_return_false: }
  }
  \cs_new_eq:NN \hep@soul@if@safe \hep_soul_if_safe:nTF
}
\ExplSyntaxOff
%    \end{macrocode}
% \end{macro}
%
% \begin{macro}{\hep@text@underline}
% Underline text with "\textul", which allows line breaks, if \software{soul} can process it.
% Otherwise, for example for text containing "\ref"\marg{key}, "\enquote"\marg{text}, or "\eg", underline it as a whole with the original "\underline".
% The macro is long, since "\textul" accepts paragraphs.
%    \begin{macrocode}
\newcommand\hep@text@underline[1]{%
  \hep@soul@if@safe{#1}{\textul{#1}}{\hep@underline{#1}}%
}
%    \end{macrocode}
% \end{macro}
%
% Use \software{pdftexcmds} \cite{pdftexcmds} to map the selected language to the option expected by \software{foreign}.
%    \begin{macrocode}
\RequirePackage{pdftexcmds}
\ifnum\pdf@strcmp{\heptext@lang}{american}=0
  \newcommand*{\heptext@lang@foreign}{USenglish}
\else
  \ifnum\pdf@strcmp{\heptext@lang}{USenglish}=0
    \newcommand*{\heptext@lang@foreign}{USenglish}
  \else
    \newcommand*{\heptext@lang@foreign}{british}
  \fi
\fi
%    \end{macrocode}
%
%\begin{macro}{\vs}
% Load \software{foreign} \cite{foreign} for common foreign abbreviations and vocabulary.
% Add the missing "\vs" abbreviation.
%    \begin{macrocode}
\PassOptionsToPackage{all, \heptext@lang@foreign}{foreign}
\RequirePackage{foreign}
\renewcommand*\foreignabbrfont{}
\DeclareRobustCommand\vs{\xperiodafter{{\foreignabbrfont{vs}}}}
%    \end{macrocode}
% \end{macro}
%
% The \software{foreign} package relies on the \software{xspace} package \cite{xspace}.
% Ensure that "\xspace" is compatible with the "\enquote" macro from \software{csquotes}.
%    \begin{macrocode}
\xspaceaddexceptions{\csq@qclose@i}
%    \end{macrocode}
%
%\begin{macro}{\no}
% Define "\no"\marg{number} for \textnumero\ with appropriate spacing.
%    \begin{macrocode}
\newcommand*{\no}[1]{\textnumero~#1}
%    \end{macrocode}
% \end{macro}
%
%\begin{macro}{\software}
% Define "\software"\linebreak[1]\oarg{version}\linebreak[1]\marg{name} for software names with optional version information, using \software{relsize} \cite{relsize} and \software{xparse} \cite{xparse}.
%    \begin{macrocode}
\RequirePackage{relsize}
\newcommand*{\software}[2][\hspace{-\fontdimen2\font}]{%
  {\smaller[.5]\textsc{#2}~#1}%
}
\RequirePackage{xparse}
\NewExpandableDocumentCommand\hep@gobble@opt@arg{om}{#2}
%    \end{macrocode}
% \end{macro}
%
% \begin{macro}{\hep@url@text}
% Print a URL or email address in typewriter font without \software{hyperref}.
% As for "\nolinkurl", the text is printed verbatim, so that characters such as "_", "#", "&", and "~" need no escaping.
% The escaped characters "\%", "\#", "\&", "\_", "\$", and "\~" become the characters themselves.
% Macros in the text, such as the fields of \software{biblatex}, are expanded first.
% The \software{hep-title} and \software{hep-bibliography} packages define the same macro, so define it only once.
%    \begin{macrocode}
\ExplSyntaxOn
\cs_if_exist:NF \hep@url@text {
  \tl_new:N \l__hep_url_tl
  \cs_generate_variant:Nn \tl_replace_all:Nnn { NnV }
  \cs_new_protected:Npn \hep@url@text #1 {
    \tl_set:Nn \l__hep_url_tl {#1}
    \regex_replace_all:nnN { \cP. } { \cO\# } \l__hep_url_tl
    \protected@edef \l__hep_url_tl { \l__hep_url_tl }
    \tl_remove_all:Nn \l__hep_url_tl { \protect }
    \tl_replace_all:NnV \l__hep_url_tl { \# } \c_hash_str
    \tl_replace_all:NnV \l__hep_url_tl { \% } \c_percent_str
    \tl_replace_all:NnV \l__hep_url_tl { \& } \c_ampersand_str
    \tl_replace_all:NnV \l__hep_url_tl { \_ } \c_underscore_str
    \tl_replace_all:NnV \l__hep_url_tl { \$ } \c_dollar_str
    \tl_replace_all:NnV \l__hep_url_tl { \~ } \c_tilde_str
    \texttt { \tl_to_str:N \l__hep_url_tl }
  }
}
\ExplSyntaxOff
%    \end{macrocode}
% \end{macro}
%
% \begin{macro}{\online}
% \begin{macro}{\email}
% Define "\online"\marg{url}\marg{text}, combining the link target of "\href" with the URL-style text of "\url".
% Define "\email"\marg{address} as a convenience wrapper.
%    \begin{macrocode}
\providecommand*{\online}[2]{\hep@url@text{#2}}%
\providecommand*{\hep@email}[1]{\online{mailto:#1}{#1}}
\providecommand*{\email}{\hep@email}
\AtBeginDocument{\@ifpackageloaded{hyperref}{%
  \renewcommand*{\online}[2]{\href{#1}{\nolinkurl{#2}}}%
}{}}
%    \end{macrocode}
% \end{macro}
% \end{macro}
%
% Make the text macros safe in \PDF strings such as bookmarks.
% Software names, underlined text, links, and email addresses reduce to their text.
% The abbreviations of \software{foreign} end with "\xperiodafter" or "\xperiodcommaafter" from \software{xpunctuate}, which look ahead and are therefore replaced by their expandable core.
% \begin{macro}{\hep@pdf@period}
% \begin{macro}{\hep@pdf@periodcomma}
% Add the period, or period and comma, followed by "\xspace", which \software{hyperref} supports in \PDF strings.
%    \begin{macrocode}
\newcommand*\hep@pdf@period[1]{#1.\xspace}
\newcommand*\hep@pdf@periodcomma[1]{#1.,\xspace}
\AtBeginDocument{\@ifpackageloaded{hyperref}{%
  \pdfstringdefDisableCommands{%
    \let\software\hep@gobble@opt@arg
    \let\underline\@firstofone
    \let\textul\@firstofone
    \let\ul\@firstofone
    \let\online\@secondoftwo
    \let\hep@email\@firstofone
    \let\xperiodafter\hep@pdf@period
    \let\xperiodcommaafter\hep@pdf@periodcomma
  }%
}{}}
%    \end{macrocode}
% \end{macro}
% \end{macro}
%
% \begin{macro}{\prefix}
% Define "\prefix"\marg{prefix}\marg{word} with a controlled line break after the prefix.
%    \begin{macrocode}
\newcommand*{\prefix}[2]{(#1\mbox{-)}\allowbreak #2}
%    \end{macrocode}
% \end{macro}
%
% \begin{macro}{\subsubparagraph}
% Define "\subsubparagraph", which starts a paragraph-level break without printing a title.
%    \begin{macrocode}
\newcommand*\subsubparagraph{%
  \@tempskipa 3.25ex \@plus 1ex \@minus .2ex\relax%
  \addpenalty\@secpenalty\addvspace\@tempskipa%
  \everypar{\setbox\z@\lastbox\everypar{}}%
}
%    \end{macrocode}
% \end{macro}
%
% \subsection{Lists}
%
% Load \software{enumitem} \cite{enumitem}.
%    \begin{macrocode}
\PassOptionsToPackage{inline}{enumitem}
\RequirePackage{enumitem}
%    \end{macrocode}
%
% \begin{environment}{inlinelist}
% Define the inline list environment.
%    \begin{macrocode}
\newlist{inlinelist}{enumerate*}{1}
\setlist*[inlinelist,1]{%
  label=\roman*), itemjoin={,\ }, itemjoin*={, and\ }, after=.%
}
%    \end{macrocode}
% \end{environment}
%
% \begin{macro}{\hep@descript@item}
% \begin{macro}{\hep@descript@select}
% \begin{environment}{enumdescript}
% Define the enumerated description list environment.
% Its "\item"\oarg{label}\marg{title} prints the title in bold, followed by a space unless it is empty.
% The redefined "\item" acts only at the list depth of the environment, so that nested lists use the original "\item".
% The original "\item" is saved only if the redefined one is not active yet, because in nested "enumdescript" lists saving the active "\item" would make it call itself.
%    \begin{macrocode}
\newlist{enumdesc}{enumerate}{2}
\setlist[enumdesc]{label=\arabic*.}
\newcommand*{\hep@descript@item}[2][]{%
  \if\relax\detokenize{#1}\relax\hep@item\else\hep@item[#1]\fi%
  \textbf{#2}\if\relax\detokenize{#2}\relax\else\hspace{1em}\fi%
  \@ifnextchar\par\@gobble\relax%
}
\newcommand*{\hep@descript@select}{%
  \ifnum\@listdepth=\hep@descript@depth\relax%
    \expandafter\hep@descript@item%
  \else%
    \expandafter\hep@item%
  \fi%
}
\newenvironment{enumdescript}[1][]{%
  \begin{enumdesc}[#1]%
  \edef\hep@descript@depth{\the\@listdepth}%
  \ifx\item\hep@descript@select\else%
    \let\hep@item\item%
    \let\item\hep@descript@select%
  \fi%
}{\end{enumdesc}}
%    \end{macrocode}
% \end{environment}
% \end{macro}
% \end{macro}
%
% Load \software{noindentafter} \cite{noindentafter} and suppress paragraph indentation after list environments.
%    \begin{macrocode}
\RequirePackage{noindentafter}
\NoIndentAfterEnv{itemize}
\NoIndentAfterEnv{enumerate}
\NoIndentAfterEnv{description}
\NoIndentAfterEnv{enumdescript}
\NoIndentAfterEnv{verbatim}
%    \end{macrocode}
%
%</package>
%
% \section{Test}
%
%<*test>
%
%    \begin{macrocode}
\documentclass{article}

\usepackage[lang=american]{hep-text}
\usepackage{hyperref}

\begin{document}

\section{HEP-Text test}

Language selection: \languagename.

\MakeOuterQuote{"}
Quotation commands: \enquote{quoted text} and "outer quoted text".

Foreign abbreviations: \eg, \ie, \cf, and \vs.
\begingroup
\renewcommand\foreignabbrfont{\itshape}
Styled foreign abbreviations: \eg, \ie, \cf, and \vs.
\endgroup

Numero labels: \no{123}.

Software names: \software{HEP-Text} and \software[v1.5]{HEP-Text}.

Online material:
\online{https://ctan.org/pkg/hep-text}
  {ctan.org/pkg/hep-text}.
Email addresses: \email{jan.hajer@tecnico.ulisboa.pt}.

Prefix constructions: \prefix{pre}{selection} and \prefix{post}{processing}.

\subsubparagraph
This paragraph follows a heading-free paragraph-level break.

\begin{itemize}
  \item An itemize entry.
\end{itemize}
This paragraph tests indentation after an itemize environment.

\begin{enumerate}
  \item An enumerate entry.
\end{enumerate}
This paragraph tests indentation after an enumerate environment.

\begin{description}
  \item[Term] A description entry.
\end{description}
This paragraph tests indentation after a description environment.

Inline lists:
\begin{inlinelist}
  \item first
  \item second
  \item third
\end{inlinelist}

\begin{enumdescript}[label=\Roman*)]
  \item{First} one
  \item{Second} two
  \item{Third} three
\end{enumdescript}
This paragraph tests indentation after an enumdescript environment.

\noindent\begin{minipage}{.35\linewidth}
\underline{This underlined sentence is long enough to wrap across
multiple lines.}
\end{minipage}

\end{document}
%    \end{macrocode}
%
%</test>
%
% \section{Readme}
%
%<*readme>
%
%    \begin{macrocode}
# The `hep-text` package

Text utilities: language setup, abbreviations, links, breakable
underlining and compact lists

## Introduction

The `hep-text` package collects lightweight text utilities. It configures
language-aware hyphenation, quotation marks, and Latin abbreviations
through established LaTeX packages; adds helpers for software names,
online references, email addresses, numero labels, prefix constructions,
and line-breakable underlining; and defines compact inline and enumerated
description lists based on `enumitem`.

Load the package with `\usepackage{hep-text}`.

## Author

Jan Hajer

## License

This file may be distributed and/or modified under the conditions of the
`LaTeX` Project Public License, either version 1.3c of this license or
(at your option) any later version. The latest version of this license is
in `http://www.latex-project.org/lppl.txt` and version 1.3c or later is
part of all distributions of LaTeX version 2005/12/01 or later.
%    \end{macrocode}
%
%</readme>
%
% \Finale

\endinput
